\section{The Schema Is the Version}
\label{sec:ch04-the-schema-is-the-version}

\index{schema!has no version number}%
Export the schema chapter~\ref{ch:data-without-n-plus-1} left behind and find
the identifier on each of the two types. The descriptions the export writes
above every field are dropped here, and the fields that are not identifiers
are elided:

\begin{graphqlsdl}
type Session {
  speaker: Speaker
  id: Int!
  ...
}

type Speaker {
  id: Int!
  ...
}
\end{graphqlsdl}

\index{identifier!collides across types}%
Session~1 and speaker~1 are different objects of different types wearing the
same identifier, because each \code{Int} is a primary key from a table and each
table starts counting at one. Inside SQLite that is fine: a key is only ever
read next to the name of the table it came from. Once it leaves through the
schema the table name is gone, and \code{1} is all the client has.

\index{client!caching by identifier}%
A client that stores objects by identifier now has to decide whether those two
are the same thing, and nothing in the response tells it.
\index{typename@\code{\_\_typename}!definition}%
The type name is available, through \code{\_\_typename}, a field every object
in every GraphQL schema answers with the name of its own type, so a careful
client can ask for it and pair the two and be safe. That is the point: the
schema has pushed a decision onto every client
that the server could have settled once, and the server is the one that knows
there are two tables.

\index{sessionById@\code{sessionById}!cannot be fixed in place}%
Now try to fix it. \code{sessionById(id: Int!)} takes an integer, so an
identifier that carries its type is a different type, a different argument and
a different field. Change it and every query a client has already written
breaks. Add a second field
beside it and the schema has two ways to fetch a session forever. There is no
third option, because there is no version number to raise. A REST service can
ship \code{/v2/sessions} and let \code{/v1} rot; a GraphQL service has one
endpoint, one schema, and every client it has ever had on it at once.

\index{schema!is the contract}%
That constraint is the whole of this chapter, and it is the deal rather than an
oversight. One endpoint means one contract, and a contract you cannot version
is a contract you have to be able to extend without breaking.
